\documentclass[11pt,a4paper]{article}

\usepackage[margin=2.4cm]{geometry}
\usepackage{amsmath,amssymb}
\usepackage{graphicx}
\usepackage{booktabs}
\usepackage{array}
\usepackage{tabularx}
\usepackage{xcolor}
\usepackage{microtype}
\usepackage{caption}
\usepackage{float}
\usepackage{listings}
\usepackage[hidelinks]{hyperref}

\definecolor{codeblue}{RGB}{32,80,130}
\definecolor{codegray}{RGB}{245,246,248}
\definecolor{coderule}{RGB}{210,214,220}

\lstdefinestyle{matlabstyle}{
  language=Matlab,
  basicstyle=\ttfamily\small,
  keywordstyle=\color{codeblue}\bfseries,
  backgroundcolor=\color{codegray},
  frame=single,
  rulecolor=\color{coderule},
  breaklines=true,
  showstringspaces=false,
  columns=fullflexible
}

\newcommand{\vect}[1]{\boldsymbol{#1}}
\newcommand{\R}{\mathbb{R}}

\title{\textbf{A Four-Level Inverse-Problem Case Study}\\
\large MATLAB implementation for a synthetic microbial consortium}
\author{Technical implementation report accompanying the manuscript\\
\emph{A Hierarchical Framework for Inverse Problems in Biological Cybernetics}}
\date{23 July 2026}

\begin{document}

\maketitle

\begin{abstract}
This report documents a reproducible MATLAB case study illustrating four
increasingly ambitious inverse problems in biological cybernetics:
parameter estimation, model discovery, inverse optimal control, and inverse
differential games. A common biological system is used throughout: two
microbial strains that compete for a supplied substrate while reciprocally
exchanging essential metabolites. The model has five dynamic states and two
metabolic-allocation controls. Synthetic experiments are generated from known
parameters and objectives, allowing the accuracy and limitations of each
inverse step to be quantified.

The six calibrated kinetic parameters are recovered with a root-mean-square
(RMS) relative error of \(0.572\%\). Sparse model discovery obtains
\(R^2=0.950\) and \(R^2=0.946\) for the two per-capita growth laws, while also
revealing a structural ambiguity for the second strain. Inverse optimal control
recovers a centralized community objective with Euclidean weight error
\(1.57\times10^{-2}\). The inverse differential-game formulation recovers the
two player-specific objectives with errors \(4.13\times10^{-2}\) and
\(6.68\times10^{-3}\), and reproduces the equilibrium controls with RMS error
\(1.02\times10^{-3}\). The complete workflow executes in approximately
13 seconds on the development system, passes all four unit tests, and produces
no MATLAB Code Analyzer findings.
\end{abstract}

\tableofcontents
\clearpage

\section{Purpose and scope}

The manuscript proposes a hierarchy in which each inverse problem inherits the
uncertainties of the levels below it and introduces a qualitatively new source
of non-uniqueness. A useful case study should therefore do more than present
four unrelated algorithms. It should hold the biological setting fixed while
changing what is assumed known and what is inferred.

The present implementation was designed around five requirements:
\begin{enumerate}
  \item use the same microbial system at all four levels;
  \item keep the dynamics small enough for transparent mathematical analysis;
  \item introduce explicit, biologically interpretable control variables for
        Levels 3 and 4;
  \item provide known ground truth through synthetic data; and
  \item expose non-identifiability rather than forcing exact recovery.
\end{enumerate}

The result is a proof-of-concept, not a claim that objectives can already be
reliably inferred from ordinary microbial abundance data. In particular,
dynamic allocation measurements or appropriate molecular proxies would be
needed before applying the upper levels to a real consortium.

\section{Biological system and dynamic model}

\subsection{Two-strain cross-feeding consortium}

Two microbial strains \(X_1\) and \(X_2\) compete for a common substrate \(S\).
Strain 1 secretes metabolite \(M_1\), which supports growth of strain 2, while
strain 2 secretes metabolite \(M_2\), which supports growth of strain 1. The
state is
\begin{equation}
  \vect{x}(t)=
  \begin{bmatrix}
    X_1(t) & X_2(t) & S(t) & M_1(t) & M_2(t)
  \end{bmatrix}^{\mathsf T}.
\end{equation}

Each strain controls a metabolic allocation \(u_i(t)\in[0,1]\). Increasing
\(u_i\) increases secretion of \(M_i\), but reduces the strain's instantaneous
growth capacity. This construction creates a biologically interpretable
trade-off between individual growth and support of the partner.

\subsection{Growth and material-balance equations}

For strain \(i\), let \(j=3-i\) denote its partner. The unburdened specific
growth rate is
\begin{equation}
  \bar{\mu}_i(\vect{x}) =
  \mu_{i,\max}
  \frac{S}{K_{S,i}+S}
  \frac{M_j}{K_{M,i}+M_j},
\end{equation}
and the allocation-dependent growth rate is
\begin{equation}
  \mu_i(\vect{x},u_i)=
  \bar{\mu}_i(\vect{x})(1-c_i u_i).
\end{equation}
The implementation enforces a small positive lower bound on the allocation
factor for numerical robustness. The complete dynamics are
\begin{align}
  \dot X_i &=
  \left(\mu_i(\vect{x},u_i)-D\right)X_i, \label{eq:biomass}\\
  \dot S &=
  D(S_{\mathrm{in}}-S)
  -\sum_{i=1}^{2}\frac{\mu_i(\vect{x},u_i)X_i}{Y_{S,i}},
  \label{eq:substrate}\\
  \dot M_1 &=
  \alpha_1u_1X_1
  -\frac{\mu_2(\vect{x},u_2)X_2}{Y_{M,2}}
  -DM_1,\\
  \dot M_2 &=
  \alpha_2u_2X_2
  -\frac{\mu_1(\vect{x},u_1)X_1}{Y_{M,1}}
  -DM_2. \label{eq:metabolites}
\end{align}

Table~\ref{tab:parameters} lists the ground-truth parameters used to generate
the synthetic data.

\begin{table}[H]
\centering
\caption{Ground-truth parameters of the synthetic microbial consortium.}
\label{tab:parameters}
\begin{tabular}{lll}
\toprule
Quantity & Strain 1 & Strain 2\\
\midrule
Maximum growth rate, \(\mu_{i,\max}\) & 0.55 & 0.50\\
Substrate half-saturation, \(K_{S,i}\) & 0.80 & 1.00\\
Metabolite half-saturation, \(K_{M,i}\) & 0.12 & 0.10\\
Allocation cost, \(c_i\) & 0.35 & 0.30\\
Secretion rate, \(\alpha_i\) & 0.55 & 0.50\\
Substrate yield, \(Y_{S,i}\) & 0.65 & 0.60\\
Metabolite yield, \(Y_{M,i}\) & 8.0 & 8.0\\
\midrule
Dilution rate, \(D\) & \multicolumn{2}{c}{0.05}\\
Inlet substrate, \(S_{\mathrm{in}}\) & \multicolumn{2}{c}{10.0}\\
\bottomrule
\end{tabular}
\end{table}

\section{Synthetic experiments and numerical implementation}

\subsection{Calibration experiments}

Three calibration experiments use different initial biomass ratios, substrate
levels, and metabolite concentrations. The controls are persistently exciting
square-wave schedules with small sinusoidal components. They span
approximately \(0.08\) to \(0.82\), which is important for separating
secretion-rate effects from allocation costs. States are sampled at 49 times
over 24 time units and corrupted with independent Gaussian noise at
\(0.5\%\) of the maximum state scale.

\subsection{Two simulation modes}

The MATLAB implementation contains two complementary integrators:
\begin{itemize}
  \item \texttt{simulateCommunity.m} uses \texttt{ode15s} with interpolated
        controls and non-negativity constraints;
  \item \texttt{simulateControlledModel.m} uses fixed-step fourth-order
        Runge--Kutta integration for the repeated simulations required by
        optimization.
\end{itemize}

The unit tests verify that the two simulators agree to an all-state RMS
difference below \(2\times10^{-3}\) on a representative experiment.

\subsection{Software entry point}

The entire workflow is launched by
\begin{lstlisting}[style=matlabstyle]
results = runInverseLadderCaseStudy;
\end{lstlisting}
The command generates data, executes the four levels, performs forward
validation, saves \texttt{inverse\_ladder\_results.mat}, and exports the
summary graphics.

\section{Level 1: parameter estimation}

\subsection{Inverse problem}

The model structure in equations~\eqref{eq:biomass}--\eqref{eq:metabolites}
is assumed known. Six quantities are estimated:
\begin{equation}
  \vect{\theta} =
  [\mu_{1,\max},\mu_{2,\max},\alpha_1,\alpha_2,c_1,c_2]^{\mathsf T}.
\end{equation}
All remaining parameters are treated as independently characterized. The
weighted least-squares problem is
\begin{equation}
  \hat{\vect{\theta}} =
  \arg\min_{\vect{\theta}_{\min}\le\vect{\theta}\le\vect{\theta}_{\max}}
  \sum_{e,k,\ell}
  \left[
    \frac{x^{\mathrm{sim}}_{e,k,\ell}(\vect{\theta})
          -x^{\mathrm{obs}}_{e,k,\ell}}
         {s_{\ell}}
  \right]^2,
\end{equation}
where \(e\), \(k\), and \(\ell\) index experiments, times, and states,
respectively, and \(s_{\ell}\) is a state-specific scale.

The problem is solved with \texttt{lsqnonlin} from two dispersed initial
points. Bound constraints enforce biologically meaningful values.

\subsection{Results}

\begin{table}[H]
\centering
\caption{Level 1 parameter-estimation results.}
\label{tab:level1}
\begin{tabular}{lrrr}
\toprule
Parameter & True & Estimated & Relative error (\%)\\
\midrule
\(\mu_{1,\max}\) & 0.55000 & 0.55080 &  0.145\\
\(\mu_{2,\max}\) & 0.50000 & 0.49919 & -0.162\\
\(\alpha_1\)     & 0.55000 & 0.54964 & -0.065\\
\(\alpha_2\)     & 0.50000 & 0.50039 &  0.078\\
\(c_1\)          & 0.35000 & 0.35175 &  0.501\\
\(c_2\)          & 0.30000 & 0.29614 & -1.287\\
\bottomrule
\end{tabular}
\end{table}

The RMS relative error across all six parameters is \(0.572\%\).
Figure~\ref{fig:level1} shows the observed and fitted trajectories for the
first calibration experiment. The fitted model reproduces biomass, substrate,
and metabolite dynamics while being driven by the measured allocation
schedules.

\begin{figure}[H]
\centering
\includegraphics[width=\textwidth]{figures/level1_calibration.png}
\caption{Level 1 calibration for the first synthetic experiment. Circles are
noisy observations; solid curves are predictions obtained with the estimated
parameters. The bottom-right panel shows the externally prescribed allocation
controls.}
\label{fig:level1}
\end{figure}

\section{Level 2: model discovery}

\subsection{Candidate library}

The environmental balances are retained, while the per-capita biomass growth
law is treated as unknown. For each strain, a sparse regression library is
constructed from:
\begin{equation}
  \vect{\phi} =
  \begin{bmatrix}
  q_i &
  q_i u_i &
  X_1 &
  X_2 &
  S &
  M_j &
  u_i
  \end{bmatrix},
  \qquad
  q_i =
  \frac{S}{K_{S,i}+S}
  \frac{M_j}{K_{M,i}+M_j}.
\end{equation}

Rather than differentiating noisy biomass measurements directly, the method
uses window-integrated per-capita growth:
\begin{equation}
  \frac{\log X_i(t_b)-\log X_i(t_a)}{t_b-t_a}
  \approx
  \beta_{i,0}
  + \frac{1}{t_b-t_a}
    \int_{t_a}^{t_b}\vect{\phi}(t)^{\mathsf T}\vect{\beta}_i\,dt.
\end{equation}
Savitzky--Golay smoothing is applied before integration. LASSO selects a sparse
coefficient vector using cross-validation.

\subsection{Results and structural ambiguity}

\begin{table}[H]
\centering
\caption{Level 2 sparse growth-law coefficients. Zeros denote terms removed
by regularization.}
\label{tab:level2}
\begin{tabular}{lrr}
\toprule
Candidate term & Strain 1 & Strain 2\\
\midrule
Intercept & -0.0529 & -0.0252\\
Monod cross-feeding, \(q_i\) & 0.5200 & 0.4157\\
Allocation interaction, \(q_i u_i\) & -0.2090 & 0\\
Strain 1 abundance & 0 & 0\\
Strain 2 abundance & 0 & 0\\
Substrate & 0 & 0\\
Received metabolite & 0 & 0\\
Direct allocation, \(u_i\) & 0 & -0.0496\\
\midrule
\(R^2\) & 0.9497 & 0.9458\\
\bottomrule
\end{tabular}
\end{table}

For strain 1, the discovered support agrees with the generating model:
growth depends on \(q_1\) and the interaction \(q_1u_1\). For strain 2, the
method selects a direct allocation term instead of \(q_2u_2\), despite
achieving \(R^2=0.946\). This is not treated as a software failure. It is an
example of structural indistinguishability caused by correlated candidate
functions over the observed region of state-control space. A high predictive
score does not guarantee recovery of the correct biological mechanism.

\begin{figure}[H]
\centering
\includegraphics[width=\textwidth]{figures/inverse_ladder_summary.png}
\caption{Summary of the four inverse levels. The Level 2 heat map displays the
selected sparse coefficients, while the other panels compare ground-truth and
estimated parameters or objective weights.}
\label{fig:summary}
\end{figure}

\section{Level 3: inverse optimal control}

\subsection{Centralized-planner hypothesis}

At Level 3, the consortium is provisionally treated as a single decision
maker choosing both allocation trajectories. This is a hypothesis to be
tested, not an assertion of community-level selection. The finite-horizon
objective is linear in four normalized kernels:
\begin{equation}
  J_C(\vect{u};\vect{w}) =
  \sum_{m=1}^{4}w_m F_m(\vect{x},\vect{u}),
  \qquad
  \vect{w}\ge 0,\qquad
  \vect{1}^{\mathsf T}\vect{w}=1.
\end{equation}
The kernels represent:
\begin{enumerate}
  \item negative total biomass, so minimization favors biomass production;
  \item quadratic allocation effort;
  \item residual substrate;
  \item squared biomass imbalance \((X_1-X_2)^2\).
\end{enumerate}

The forward optimal-control problem uses piecewise-constant controls on eight
intervals and fourth-order Runge--Kutta integration. It is solved with
\texttt{fmincon} and box constraints \(0.02\le u_i\le0.95\).

\subsection{Inverse stationarity formulation}

Because the objective is linear in \(\vect{w}\), its control gradient at an
observed optimum \(\vect{u}^{*}\) can be written
\begin{equation}
  \nabla_{\vect{u}}J_C(\vect{u}^{*};\vect{w})
  =
  \begin{bmatrix}
    \nabla_{\vect{u}}F_1 &
    \cdots &
    \nabla_{\vect{u}}F_4
  \end{bmatrix}
  \vect{w}
  = A_C\vect{w}.
\end{equation}
Finite differences provide the feature gradients. Rows associated with active
control bounds are excluded when enough interior controls are available. The
inverse problem is
\begin{equation}
  \hat{\vect{w}} =
  \arg\min_{\vect{w}\ge0,\ \vect{1}^{\mathsf T}\vect{w}=1}
  \|A_C\vect{w}\|_2^2
  +\lambda\|\vect{w}-\vect{w}_0\|_2^2,
\end{equation}
with a very small regularization parameter
\(\lambda=10^{-8}\) and uniform reference \(\vect{w}_0\).

\subsection{Results}

\begin{table}[H]
\centering
\caption{Level 3 centralized objective weights.}
\label{tab:level3}
\begin{tabular}{lrr}
\toprule
Objective kernel & True & Inferred\\
\midrule
Total biomass & 0.6200 & 0.6098\\
Control effort & 0.1400 & 0.1391\\
Residual substrate & 0.1400 & 0.1519\\
Community imbalance & 0.1000 & 0.0991\\
\bottomrule
\end{tabular}
\end{table}

The Euclidean weight error is \(1.57\times10^{-2}\). Forward solution using
the inferred weights returns the same discretized controls and states to the
reported numerical precision. The singular values of the stacked stationarity
matrix are approximately
\[
  8.23\times10^{-1},\quad
  2.31\times10^{-1},\quad
  2.2\times10^{-3},\quad
  \text{near zero}.
\]
The separation between the two smallest singular values is approximately
641, supporting a well-defined one-dimensional near-null direction after
normalization.

\begin{figure}[H]
\centering
\includegraphics[width=\textwidth]{figures/level3_ioc.png}
\caption{Level 3 results for the first experiment. The inferred objective
reproduces the observed optimal controls and biomass trajectories. The
bottom-right panel shows the singular-value spectrum of the inverse
stationarity matrix.}
\label{fig:level3}
\end{figure}

\section{Level 4: inverse differential game}

\subsection{Player-specific objectives}

The centralized hypothesis is replaced by two independent decision makers.
Player \(i\) minimizes
\begin{equation}
  J_i(u_i,u_j;\vect{w}_i)
  =
  \sum_{m=1}^{4}w_{i,m}F_{i,m},
  \qquad
  \vect{w}_i\ge0,\qquad
  \vect{1}^{\mathsf T}\vect{w}_i=1.
\end{equation}
The four kernels are:
\begin{enumerate}
  \item negative own biomass;
  \item own quadratic allocation effort;
  \item deficit of the metabolite received from the partner;
  \item negative partner biomass.
\end{enumerate}
The final term permits a general-sum interaction: supporting the partner can
be beneficial because the partner supplies an essential metabolite.

An open-loop Nash equilibrium
\((u_1^{*},u_2^{*})\) satisfies
\begin{equation}
  J_i(u_i^{*},u_j^{*})
  \le
  J_i(u_i,u_j^{*})
  \quad
  \text{for every feasible unilateral deviation }u_i.
\end{equation}

\subsection{Forward and inverse algorithms}

The forward game is solved by damped iterative best response:
\begin{enumerate}
  \item fix \(u_2\) and optimize \(u_1\);
  \item fix the updated \(u_1\) and optimize \(u_2\);
  \item damp the updates and repeat;
  \item verify that neither player obtains a material unilateral improvement.
\end{enumerate}
Each best response is an \texttt{fmincon} problem. The inverse problem is
formed separately for each player by differentiating its four feature
functionals with respect to its own control:
\begin{equation}
  \hat{\vect{w}}_i =
  \arg\min_{\vect{w}_i\ge0,\
             \vect{1}^{\mathsf T}\vect{w}_i=1}
  \|A_i\vect{w}_i\|_2^2
  +\lambda\|\vect{w}_i-\vect{w}_0\|_2^2.
\end{equation}

\subsection{Results}

\begin{table}[H]
\centering
\caption{Level 4 player-specific objective weights.}
\label{tab:level4}
\begin{tabular}{lrrrr}
\toprule
& \multicolumn{2}{c}{Player 1} & \multicolumn{2}{c}{Player 2}\\
\cmidrule(lr){2-3}\cmidrule(lr){4-5}
Objective kernel & True & Inferred & True & Inferred\\
\midrule
Own biomass & 0.5800 & 0.5679 & 0.6200 & 0.6148\\
Control effort & 0.2200 & 0.2110 & 0.2500 & 0.2506\\
Metabolite deficit & 0.0500 & 0.0856 & 0.0500 & 0.0541\\
Partner biomass & 0.1500 & 0.1355 & 0.0800 & 0.0804\\
\bottomrule
\end{tabular}
\end{table}

The Euclidean weight errors are \(4.13\times10^{-2}\) and
\(6.68\times10^{-3}\) for Players 1 and 2. Forward validation gives:
\begin{itemize}
  \item mean control RMS error \(1.02\times10^{-3}\);
  \item mean state RMS error \(9.32\times10^{-4}\);
  \item maximum post-solution unilateral improvement below
        \(7.4\times10^{-6}\).
\end{itemize}
The near-null-direction separation ratios are approximately 20.2 and 64.0 for
the two players. Player 1 is less well-conditioned, consistent with its larger
weight error, particularly for the received-metabolite term.

\begin{figure}[H]
\centering
\includegraphics[width=\textwidth]{figures/level4_idg.png}
\caption{Level 4 results for the first experiment. Dashed forward-validation
controls and trajectories closely track the equilibrium data. The lower
panels compare true and inferred objectives and show the two stationarity
spectra.}
\label{fig:level4}
\end{figure}

\section{Implementation map and verification}

\begin{table}[H]
\centering
\caption{Principal MATLAB components.}
\label{tab:software}
\begin{tabularx}{\textwidth}{>{\ttfamily}l X}
\toprule
\normalfont File & Purpose\\
\midrule
runInverseLadderCaseStudy.m & Top-level reproducible workflow.\\
defaultInverseLadderConfig.m & Parameters, experiments, objectives, and solver settings.\\
communityRhs.m & Five-state controlled microbial ODE.\\
simulateCommunity.m & Adaptive \texttt{ode15s} simulator.\\
simulateControlledModel.m & Fixed-step RK4 simulator used in optimization.\\
runLevel1ParameterEstimation.m & Bound-constrained multi-experiment calibration.\\
runLevel2ModelDiscovery.m & Integrated sparse growth-law discovery.\\
solveCommunityPlanner.m & Centralized forward optimal-control solver.\\
runLevel3InverseOptimalControl.m & Level 3 inverse stationarity and validation.\\
solveOpenLoopNash.m & Damped iterative best-response solver.\\
runLevel4InverseDifferentialGame.m & Player-wise inverse stationarity and Nash validation.\\
testInverseLadderCaseStudy.m & MATLAB unit tests.\\
\bottomrule
\end{tabularx}
\end{table}

The verified development environment is MATLAB R2026a Update 1 with
Optimization Toolbox, Global Optimization Toolbox, Statistics and Machine
Learning Toolbox, Signal Processing Toolbox, Symbolic Math Toolbox, Parallel
Computing Toolbox, and Control System Toolbox. The current implementation uses
the optimization, statistics, and signal-processing functionality directly.

\begin{table}[H]
\centering
\caption{Verification summary.}
\label{tab:verification}
\begin{tabular}{lr}
\toprule
Check & Result\\
\midrule
Complete workflow runtime without plotting & approximately 13 s\\
MATLAB unit tests & 4 of 4 passed\\
MATLAB Code Analyzer findings & 0\\
Saved numerical result file & verified\\
Saved summary figure & verified\\
Nash unilateral-improvement check & \(<7.4\times10^{-6}\)\\
\bottomrule
\end{tabular}
\end{table}

\section{Interpretation and limitations}

\subsection{What the case study demonstrates}

The case study provides a computationally coherent realization of the inverse
ladder:
\begin{itemize}
  \item Level 1 recovers continuous kinetic parameters when the structure and
        controls are known.
  \item Level 2 admits competing structural explanations even when prediction
        is accurate.
  \item Level 3 recovers an objective only after choosing candidate kernels,
        normalizing their scale, and assuming centralized optimality.
  \item Level 4 requires an equilibrium concept, player-specific actions, and
        separate stationarity conditions.
\end{itemize}

The increasing data and modeling requirements are visible in the code. Levels
1 and 2 use ordinary time-series observations and prescribed inputs. Levels 3
and 4 additionally require controls, a candidate objective basis, and a
forward optimal-control or game solver.

\subsection{Important limitations}

\begin{enumerate}
  \item \textbf{Synthetic ground truth.} The trajectories, objectives, and
        controls are generated by the same model class used for inference.
        This is appropriate for verification but does not establish biological
        validity.
  \item \textbf{Observed controls.} Real microbial studies rarely measure
        metabolic allocations directly. Transcript, protein, or secretion
        measurements would be required as proxies.
  \item \textbf{Objective basis dependence.} Objective recovery is conditional
        on the selected feature library. An omitted kernel can redistribute
        inferred weights without visibly degrading the fit.
  \item \textbf{Open-loop equilibrium.} Level 4 infers finite-horizon open-loop
        strategies. General nonlinear feedback Nash equilibria would require
        solving coupled Hamilton--Jacobi equations and are not attempted.
  \item \textbf{Two-player scale.} The best-response and finite-difference
        methods are suitable for a compact illustration, not yet for large
        microbial communities.
  \item \textbf{Local optimality information.} Inverse stationarity establishes
        compatibility with first-order conditions. Forward re-optimization and
        unilateral-deviation tests are therefore essential.
\end{enumerate}

\section{Recommended extensions}

The most informative next developments are:
\begin{enumerate}
  \item design additional asymmetric perturbations specifically to resolve
        the Level 2 allocation ambiguity;
  \item add measurement noise to the Level 3 and Level 4 controls and quantify
        confidence regions for the inferred weights;
  \item compare the centralized and game formulations on the same behavior to
        demonstrate cooperative mimicry explicitly;
  \item replace synthetic allocations with gene-expression or metabolite-flux
        measurements from an engineered two-strain consortium;
  \item add profile-based or Bayesian uncertainty quantification for objective
        weights;
  \item test whether inferred objectives predict behavior under dilution,
        nutrient, or partner-knockdown conditions excluded from inference.
\end{enumerate}

\section{Conclusions}

The MATLAB implementation shows that all four levels can be represented in a
single biological and computational framework. Parameter estimation and sparse
model discovery are inexpensive and mature enough for routine application.
Inverse optimal control is feasible for a small controlled community when
allocation trajectories are available. A two-player inverse differential game
is also computationally feasible under open-loop, finite-horizon assumptions,
but it introduces a substantially larger modeling commitment and more severe
conditioning issues.

The central lesson is not that the upper levels produce uniquely correct
biological objectives. Rather, the implementation demonstrates how to formulate
those hypotheses, diagnose their identifiability, and test their predictive
consequences. The Level 2 structural ambiguity and the Level 4 conditioning
difference between players are therefore substantive results: they make the
case study consistent with the manuscript's emphasis on both opportunities and
limitations.

\begin{thebibliography}{9}

\bibitem{BangaVillaverde2026}
J. R. Banga and A. F. Villaverde,
\emph{A Hierarchical Framework for Inverse Problems in Biological
Cybernetics: Opportunities and Limitations},
manuscript, 2026.

\bibitem{BruntonKutz2022}
S. L. Brunton and J. N. Kutz,
\emph{Data-Driven Science and Engineering: Machine Learning, Dynamical
Systems, and Control},
Cambridge University Press, 2022.

\bibitem{Molloy2022}
T. L. Molloy, J. Inga, S. Hohmann, and T. Perez,
\emph{Inverse Optimal Control and Inverse Noncooperative Dynamic Game Theory},
Springer, 2022.

\bibitem{Zomorrodi2012}
A. R. Zomorrodi and C. D. Maranas,
``OptCom: A multi-level optimization framework for the metabolic modeling and
analysis of microbial communities,''
\emph{PLoS Computational Biology}, vol. 8, no. 2, e1002363, 2012.

\end{thebibliography}

\end{document}
